% latex/supplement/experiments.tex
%
% Descriptions of the ten experiments. For each experiment the
% supplement records the purpose, the inputs, the artefact paths, and
% the thesis table the results feed. A reviewer running the pipeline
% should be able to identify every output directory mentioned here
% by name.

\section{Experiment descriptions}

This section documents the ten experiments the A3E-DNS evaluation
reports. Each experiment is a standalone Python script under
\texttt{experiments\_v2/}, and each is invoked by the orchestrator in
\texttt{extended\_experimentation/experiments/e1\_e10\_orchestrator.py}.
The descriptions below follow the order the orchestrator runs them in.

\subsection{E1 --- Attack scenario catalogue}

\textbf{Purpose.} Produce the ten-scenario attack catalogue with
three severity variants per scenario and three protocols per scenario,
plus a benign baseline.

\textbf{Inputs.} None; the catalogue is fully synthetic, parameterised
by the seed.

\textbf{Outputs.}
\begin{itemize}
\item \texttt{runs/experiments/E1/scenario\_sweep.json} — per-scenario
      accuracy, FPR, precision, recall.
\item \texttt{runs/experiments/E1/predictions.jsonl} — per-window
      predictions.
\item \texttt{runs/experiments/E1/manifest.json} — the run manifest.
\end{itemize}

\textbf{Feeds.} Table 30 (scenario breakdown), Table 38 (adversarial
robustness), Table 39 (statistical significance).

\subsection{E2 --- Byzantine robustness sweep}

\textbf{Purpose.} Sweep the Byzantine fraction from $5\%$ to $40\%$
across the eight attack strategies and the four schedules.

\textbf{Inputs.} The E1 predictions (as the honest baseline).

\textbf{Outputs.}
\begin{itemize}
\item \texttt{runs/experiments/E2/byzantine\_sweep.json} — per-round
      attack ratios and aggregate accuracies.
\item \texttt{runs/experiments/E2/manifest.json}.
\end{itemize}

\textbf{Feeds.} Table 31 (Byzantine robustness). The
\texttt{attack\_to\_reference\_ratio} at each fraction is what the
robustness figure plots.

\subsection{E3 --- Differential-privacy sweep}

\textbf{Purpose.} Sweep the target $\eps$ from $0.1$ to $1.0$ and
measure the utility loss at each value.

\textbf{Inputs.} The E1 predictions.

\textbf{Outputs.}
\begin{itemize}
\item \texttt{runs/experiments/E3/dp\_sweep.json} — per-$\eps$
      accuracy and FPR.
\item \texttt{runs/experiments/E3/manifest.json}.
\end{itemize}

\textbf{Feeds.} Table 32 (DP epsilon sweep). The reported
$\eps = 0.68$ is the point at which the utility loss first crosses
$0.5\%$, which the sweep identifies from the data rather than from a
pre-set threshold.

\subsection{E4a--E4d --- QHE primitives}

\textbf{Purpose.} Four separate benchmarks of the homomorphic
encryption primitives:

\begin{itemize}
\item E4a: negacyclic NTT vs.\ naive polynomial multiplication;
\item E4b: CDT Gaussian sampler correctness and timing;
\item E4c: selective homomorphic operations over the informative
      feature subset;
\item E4d: packed ciphertext size and compression ratio.
\end{itemize}

\textbf{Inputs.} All four are self-contained; none depends on the
attack catalogue.

\textbf{Outputs.} Four JSON manifests under
\texttt{runs/experiments/E4a/} through \texttt{E4d/}.

\textbf{Feeds.} Tables 33 and 34.

\subsection{E5 --- End-to-end P99 latency}

\textbf{Purpose.} Measure the P99 latency of the full pipeline under
combined attack load.

\textbf{Inputs.} The E2 and E3 artefacts, replayed through the
pipeline.

\textbf{Outputs.}
\begin{itemize}
\item \texttt{runs/experiments/E5/latency.json} — per-protocol P50,
      P95, P99, P99.9.
\item \texttt{runs/experiments/E5/manifest.json}.
\end{itemize}

\textbf{Feeds.} Table 35 (P99 latency). The reported
$1.42\,\mathrm{ms}$ is the P99 over the DoH, DoT and DoQ streams
pooled; the per-protocol values are in the same table as a
footnote.

\subsection{E6 --- Peak throughput}

\textbf{Purpose.} Measure peak QPS at concurrency levels from 1 to
128.

\textbf{Inputs.} The E5 pipeline configuration.

\textbf{Outputs.}
\begin{itemize}
\item \texttt{runs/experiments/E6/throughput.json} — per-concurrency
      QPS.
\item \texttt{runs/experiments/E6/manifest.json}.
\end{itemize}

\textbf{Feeds.} Table 35 (peak throughput). The reported
$22.5\,\mathrm{M}\,\mathrm{QPS}$ is the maximum over the sweep,
measured at concurrency $= 32$.

\subsection{E7 --- Component ablation}

\textbf{Purpose.} Ablate each of AEE, QHE, FL and the attention block
in turn, and measure the resulting accuracy, FPR and P99.

\textbf{Inputs.} Pre-computed checkpoints for the four ablated
variants, produced by re-running the training pipeline with the
corresponding component disabled.

\textbf{Outputs.}
\begin{itemize}
\item \texttt{runs/experiments/E7/ablation.json}.
\item \texttt{runs/experiments/E7/manifest.json}.
\end{itemize}

\textbf{Feeds.} Table 36. The ablation is \emph{partial}: a full
factorial design was not performed, as Section~6.7.3 of the thesis
states.

\subsection{E8 --- Multi-protocol accuracy}

\textbf{Purpose.} Confirm the detector generalises across DoH, DoT
and DoQ.

\textbf{Inputs.} The E1 catalogue, stratified by protocol.

\textbf{Outputs.}
\begin{itemize}
\item \texttt{runs/experiments/E8/multi\_protocol.json}.
\item \texttt{runs/experiments/E8/manifest.json}.
\end{itemize}

\textbf{Feeds.} Table 37.

\subsection{E9 --- Adversarial robustness}

\textbf{Purpose.} Measure the accuracy drop under the adaptive-blend
variant of each scenario, relative to the baseline severity.

\textbf{Inputs.} The E1 and E2 predictions.

\textbf{Outputs.}
\begin{itemize}
\item \texttt{runs/experiments/E9/adversarial\_robustness.json}.
\item \texttt{runs/experiments/E9/manifest.json}.
\end{itemize}

\textbf{Feeds.} Table 38. The reported $96.5\%$ robustness is the
ratio of adaptive-blend accuracy to baseline accuracy, not the raw
accuracy itself.

\subsection{E10 --- Statistical significance}

\textbf{Purpose.} Run the four significance tests (paired bootstrap,
Wilcoxon signed-rank, permutation, Friedman) over the E1--E9
artefacts.

\textbf{Inputs.} The E1, E2, E3, E5, E6, E7, E8 and E9 manifests and
prediction files.

\textbf{Outputs.}
\begin{itemize}
\item \texttt{runs/experiments/E10/statistical\_analysis.json}.
\item \texttt{runs/experiments/E10/significance\_table.tex}.
\end{itemize}

\textbf{Feeds.} Table 39. E10 depends on every other experiment, so
it is the last node in the orchestrator's DAG.
